\documentclass[11pt]{article}
\usepackage[margin=1in]{geometry}
\usepackage{amsmath,amssymb,amsthm}
\usepackage{xurl}
\usepackage{hyperref}
\sloppy
\let\oldus\_ \renewcommand{\_}{\oldus\allowbreak}
\newtheorem{theorem}{Theorem}
\newtheorem{lemma}[theorem]{Lemma}
\theoremstyle{definition}
\newtheorem{definition}[theorem]{Definition}
\newtheorem{remark}[theorem]{Remark}
\newcommand{\PSL}{\mathrm{PSL}}
\newcommand{\SL}{\mathrm{SL}}
\newcommand{\PG}{\mathrm{PG}}
\newcommand{\Stab}{\mathrm{Stab}}

\title{Proof of conjecture 00000001012:\\ the subdegrees of $\PSL(n,q)$ on $\PG(n-1,q)$ are $1$ and $(q^n-q)/(q-1)$}
\author{Jackmeson1}
\date{2026-10-05}

\begin{document}
\maketitle

\section{The conjecture and its reading}

\textbf{Statement (English, verbatim).} Definition: The permutation representation of minimal
degree of classical groups (acting on the points of projective space). Conjecture: The subdegrees
of $\PSL(n,q)$ acting on $\PG(n-1,q)$ are of the explicit $(q^n-q)/(q-1)$ type (a subdegree
formula). (The Chinese text states the same.)

\textbf{Reading.} The subdegrees of a transitive permutation group are the lengths of the orbits
of a point stabilizer. Every transitive action has the trivial subdegree $1$ (the orbit $\{p\}$ of
the fixed point $p$ itself), so ``the subdegrees are of the $(q^n-q)/(q-1)$ type'' is read as: the
subdegrees are $1$ and $(q^n-q)/(q-1)$, i.e.\ the only nontrivial subdegree is $(q^n-q)/(q-1)$.
We prove the complete statement: for every point $p$ the stabilizer has \emph{exactly two}
orbits, $\{p\}$ and the set of all other points, of sizes $1$ and $(q^n-q)/(q-1)$. This covers both
the ``full multiset of subdegrees is $[1,(q^n-q)/(q-1)]$'' reading and the ``nontrivial subdegree
equals $(q^n-q)/(q-1)$'' reading. The reading in which \emph{every} subdegree, including the
trivial one, equals $(q^n-q)/(q-1)$ is false for every transitive action on more than one point
(the orbit $\{p\}$ has size $1\neq(q^n-q)/(q-1)$, which is at least $2$), and we do not adopt it.

\textbf{Scope.} We take $n\ge 2$ (for $n=1$, $\PG(0,q)$ is a single point, which we also prove,
and there is no nontrivial subdegree) and $q$ the order of an arbitrary finite field $F$. The
Definition line's phrase ``of minimal degree'' is not used: the conjecture names the action of
$\PSL(n,q)$ on the points of $\PG(n-1,q)$ explicitly, and that is the action we treat. We do not
prove that this action has minimal degree among faithful transitive actions.

\section{Definitions (as formalized)}

\begin{definition}
Let $F$ be a finite field with $q=|F|$ elements and $n\ge 0$. $\SL(n,F)$ is the group of
$n\times n$ matrices over $F$ of determinant $1$ (Mathlib \texttt{Matrix.SpecialLinearGroup (Fin n) F}).
$\PSL(n,F)=\SL(n,F)/Z(\SL(n,F))$, the quotient by the centre (Mathlib
\texttt{Matrix.ProjectiveSpecialLinearGroup (Fin n) F}; Lean abbreviation \texttt{PSL n F}).
\end{definition}

\begin{definition}
The point set of $\PG(n-1,F)$ is the set of one-dimensional subspaces of $F^n$, i.e.\ nonzero
vectors of $F^n$ modulo nonzero scalars (Mathlib \texttt{Projectivization F (Fin n -> F)}; Lean
abbreviation \texttt{PG n F}). $\SL(n,F)$ acts by $g\cdot[v]=[gv]$, and $\PSL(n,F)$ acts through
$\SL(n,F)$: Mathlib defines the action of $\PSL$ as the factorization of the $\SL$ permutation
representation through $\SL/Z$, and \texttt{Matrix.ProjectiveSpecialLinearGroup.smul\_proj\_mk}
states $\bar g\cdot x=g\cdot x$ for the image $\bar g$ of $g\in\SL(n,F)$ (our \texttt{psl\_smul\_eq}).
\end{definition}

\begin{definition}
For a point $p$, $\Stab(p)=\{g\in\PSL(n,F): g\cdot p=p\}$ (Mathlib \texttt{MulAction.stabilizer}),
acting on the points by restriction. The orbit of $x$ under $\Stab(p)$ is
$\Stab(p)\cdot x=\{h\cdot x : h\in\Stab(p)\}$ (Mathlib \texttt{MulAction.orbit}); the set of
orbits is the quotient \texttt{MulAction.orbitRel.Quotient}. The \emph{subdegrees} are the
cardinalities \texttt{Nat.card} of these orbits.
\end{definition}

\textbf{Conventions.} Natural-number arithmetic: $q^n-q$ and $q-1$ are truncated differences and
$/$ is floor division in $\mathbb{N}$; we prove that the division is exact,
$(q-1)\cdot\frac{q^n-q}{q-1}=q^n-q$, so the formula has its usual meaning.

\section{Proof}

\begin{lemma}[2-transitivity; \texttt{psl\_two\_transitive}]
If $a\neq b$ and $c\neq d$ are points of $\PG(n-1,F)$, there is $g\in\PSL(n,F)$ with
$g\cdot a=c$ and $g\cdot b=d$.
\end{lemma}
\begin{proof}
Mathlib proves that $\SL(n,F)$ acts 2-transitively on the points
(instance \texttt{IsMultiplyPretransitive (Matrix.SpecialLinearGroup (Fin n) F) (P F (Fin n -> F)) 2},
derived from \texttt{Projectivization.specialLinearGroup\_is\_two\_pretransitive}, file
\texttt{Mathlib/LinearAlgebra/Projectivization/Action.lean}, by D.~Loeffler and A.~Chambert-Loir):
a linear isomorphism carrying representatives of $a,b$ to representatives of $c,d$ exists because
two distinct points have linearly independent representatives, and its determinant is corrected
to $1$ by a dilatation fixing $a$ and $b$. By \texttt{MulAction.is\_two\_pretransitive\_iff} this gives
$g\in\SL(n,F)$ with $g\cdot a=c$, $g\cdot b=d$, and its image $\bar g\in\PSL(n,F)$ acts the same way.
\end{proof}

\begin{lemma}[\texttt{orbit\_stabilizer\_self}, \texttt{orbit\_stabilizer\_of\_ne}]
$\Stab(p)\cdot p=\{p\}$, and for $x\ne p$, $\Stab(p)\cdot x=\PG(n-1,F)\setminus\{p\}$.
\end{lemma}
\begin{proof}
The first is the definition of the stabilizer. For $x\ne p$: if $h\in\Stab(p)$ and $h\cdot x=p=h\cdot p$
then $x=p$ by injectivity of $h$, so $\Stab(p)\cdot x\subseteq\{p\}^c$. Conversely, for $y\ne p$
apply the previous lemma to $(p,x)\mapsto(p,y)$: the element $g$ found fixes $p$ and sends $x$ to $y$.
\end{proof}

\begin{lemma}[point count; \texttt{pow\_eq\_card\_points}, \texttt{card\_points}, \texttt{arith}]
Let $N$ be the number of points. Then $q^n=N(q-1)+1$, i.e.\ $N=(q^n-1)/(q-1)$. If $n\ge 2$ then
$N\ge 3$, $N-1=(q^n-q)/(q-1)$ and $(q-1)\cdot\frac{q^n-q}{q-1}=q^n-q$.
\end{lemma}
\begin{proof}
The count is Mathlib's \texttt{Projectivization.card'} and \texttt{Projectivization.card''}
(nonzero vectors are partitioned into $N$ classes of size $q-1$), with $|F^n|=q^n$.
Write $q=r+2$. Since $q^n\ge q^2$, $N(r+1)+1\ge (r+2)^2$ forces $N\ge 3$. Writing $N=M+1$,
$q^n-q=M(r+1)+(r+2)-(r+2)=M(q-1)$, so $(q^n-q)/(q-1)=M=N-1$, with exact division.
\end{proof}

\begin{theorem}[main theorem, \texttt{Conjecture1012.subdegrees\_PSL}]
Let $F$ be a finite field with $q$ elements, $n\ge 2$, and $p$ any point of $\PG(n-1,F)$. Then
\begin{enumerate}
\item $\PSL(n,F)$ acts transitively and faithfully on the points of $\PG(n-1,F)$;
\item the set of orbits of $\Stab(p)$ is exactly $\{\{p\},\ \{p\}^c\}$, these two sets are
  different, and the number of orbits (\texttt{Nat.card} of the orbit quotient) is $2$;
\item $|\Stab(p)\cdot p|=1$ and $|\Stab(p)\cdot x|=(q^n-q)/(q-1)$ for every $x\neq p$;
\item $(q-1)\cdot\frac{q^n-q}{q-1}=q^n-q$ and $\frac{q^n-q}{q-1}=\frac{q^n-1}{q-1}-1$;
\item the set of subdegrees $\{|\Stab(p)\cdot x| : x\}$ equals $\{1,(q^n-q)/(q-1)\}$, and
  $1\neq(q^n-q)/(q-1)$.
\end{enumerate}
\end{theorem}
\begin{proof}
Transitivity and faithfulness are found by instance search in Mathlib (Mathlib proves the action primitive, which includes transitivity;
\texttt{FaithfulSMul} holds since the kernel of the $\SL$ action is the centre,
\texttt{Projectivization.SL\_mulAction\_ker}). Since $N\ge 3$ there is a point $x_0\neq p$, so
$\{p\}\neq\{p\}^c$. Item 2 follows from the orbit lemma, and the orbit quotient injects into sets
of points (\texttt{orbitRel.Quotient.orbit\_injective}) with image $\{\{p\},\{p\}^c\}$, so it has
$2$ elements. Item 3: $|\{p\}|=1$ and $|\{p\}^c|=N-1=(q^n-q)/(q-1)$. Item 4 is the counting lemma.
Item 5 follows from items 2 and 3, and $(q^n-q)/(q-1)=N-1\ge 2$.
\end{proof}

\begin{remark}[\texttt{card\_points\_one}]
For $n=1$ the space $\PG(0,F)$ has exactly one point, so the hypothesis $n\ge 2$ is needed for a
nontrivial subdegree to exist.
\end{remark}

\section{Lean formalization and axiom audit}

File \texttt{lean/Conjecture1012/Basic.lean} (namespace \texttt{Conjecture1012}, only
\texttt{import Mathlib}). Main theorem:
{\small
\begin{verbatim}
theorem subdegrees_PSL (hn : 2 <= n) (p : PG n F) :
    IsPretransitive (PSL n F) (PG n F) /\ FaithfulSMul (PSL n F) (PG n F) /\
    Set.range (fun x : PG n F => orbit (stabilizer (PSL n F) p) x) =
      {{p}, {p}^c} /\
    ({p} : Set (PG n F)) != {p}^c /\
    Nat.card (orbitRel.Quotient (stabilizer (PSL n F) p) (PG n F)) = 2 /\
    Nat.card (orbit (stabilizer (PSL n F) p) p) = 1 /\
    (forall x : PG n F, x != p -> Nat.card (orbit (stabilizer (PSL n F) p) x) =
      (Fintype.card F ^ n - Fintype.card F) / (Fintype.card F - 1)) /\
    (Fintype.card F - 1) *
      ((Fintype.card F ^ n - Fintype.card F) / (Fintype.card F - 1)) =
      Fintype.card F ^ n - Fintype.card F /\
    (Fintype.card F ^ n - Fintype.card F) / (Fintype.card F - 1) =
      (Fintype.card F ^ n - 1) / (Fintype.card F - 1) - 1 /\
    Set.range (fun x : PG n F => Nat.card (orbit (stabilizer (PSL n F) p) x)) =
      {1, (Fintype.card F ^ n - Fintype.card F) / (Fintype.card F - 1)} /\
    1 != (Fintype.card F ^ n - Fintype.card F) / (Fintype.card F - 1)
\end{verbatim}
}
(ASCII rendering; the source uses the Unicode symbols.) Here $F$ is any type with
\texttt{[Field F] [Fintype F]}, \texttt{PSL n F} is
\texttt{Matrix.ProjectiveSpecialLinearGroup (Fin n) F} and \texttt{PG n F} is
\texttt{Projectivization F (Fin n -> F)}.
Supporting theorems: \texttt{psl\_smul\_eq}, \texttt{psl\_two\_transitive},
\texttt{orbit\_stabilizer\_self}, \texttt{orbit\_stabilizer\_of\_ne}, \texttt{pow\_eq\_card\_points},
\texttt{arith}, \texttt{card\_points}, \texttt{card\_compl\_point}, \texttt{subdegree\_ge\_two},
\texttt{card\_points\_one}.

\textbf{Axiom audit.} \texttt{lake build} succeeds, and
\texttt{\#print axioms Conjecture1012.subdegrees\_PSL} and
\texttt{\#print axioms Conjecture1012.card\_points\_one} both report
\texttt{[propext, Classical.choice, Quot.sound]}. There is no \texttt{sorry}, no
\texttt{native\_decide} and no added axiom.

\textbf{Credit.} The 2-transitivity of $\SL$ on projective space, the $\PSL$ action, its
faithfulness and the point count are Mathlib results
(\texttt{Mathlib/LinearAlgebra/Projectivization/Action.lean} and \texttt{Cardinality.lean}).

\end{document}
